\documentclass[10pt,letterpaper]{article}
\usepackage[margin=1in]{geometry}
\usepackage{amsmath,amssymb,amsfonts}
\usepackage{booktabs}
\usepackage{multirow}
\usepackage{graphicx}
\usepackage{xcolor}
\usepackage{hyperref}
\usepackage{natbib}
\usepackage{enumitem}
\usepackage{array}
\usepackage{caption}
\usepackage{subcaption}
\usepackage{tabularx}

\newcolumntype{C}{>{\centering\arraybackslash}X}
\newcolumntype{L}{>{\raggedright\arraybackslash}X}

\title{Diagnostic Analysis of Value Overestimation from Temporal Abstraction in DreamerV3}

\author{
Anonymous Authors
}

\date{}

\begin{document}
\maketitle

%% ============================================================
%%  ABSTRACT
%% ============================================================
\begin{abstract}
Temporal abstraction---acting at variable time scales---is a long-sought capability in reinforcement learning, yet integrating it into world-model agents remains poorly understood.
We present a systematic diagnostic analysis of variable-rate imagination in DreamerV3, identifying \emph{Conditioning-Induced Value Overestimation} (CIVO) as a fundamental failure mode.
Through controlled factorial experiments across four environments (Walker Walk, Cheetah Run, Cartpole Swingup, Hopper Hop), we establish a three-link causal chain: (1)~a \emph{trained} temporal embedding learns dt-specific representations on RSSM states whose fidelity degrades at $\mathrm{dt}>1$ (reward prediction gap $2.78\times$); (2)~the critic bootstraps from these imprecise-but-confidently-conditioned imagined states, producing pathological value overestimation (Normalized Value Excess $\geq 1.5\times$ in 9/9 trained-embedding seeds, 0/12 untrained; AUROC $= 0.969$); (3)~the actor exploits overestimated value regions, triggering an irreversible positive-feedback collapse.
A $2\times 2$ ablation matrix confirms that \emph{both} embedding training and exposure to $\mathrm{dt}>1$ data are necessary conditions for CIVO.
We further evaluate a systematic \emph{Intervention Hierarchy} spanning gradient-level (Selective Gradient Stopping, 1/17 seeds), output-level (Reality-Anchored Value Estimation, 0/2), and representation-level (Temporal Conditioning Noise, 2/8) interventions across 42 seeds, 8 methods, and 3 environments.
No post-hoc intervention reliably eliminates CIVO, establishing that architectural solutions---such as hierarchical temporal abstraction---are necessary for safe variable-rate imagination in world-model RL.
We additionally identify that gradient magnitude is decoupled from causal importance in multi-head world models (reward head: 38--49\% of gradients but causally irrelevant; observation decoder: 16--22\% but causally critical), a methodological insight with implications beyond this work.
\end{abstract}

%% ============================================================
%%  1  INTRODUCTION
%% ============================================================
\section{Introduction}

Temporal abstraction---the ability to reason and act at multiple time scales---is a foundational goal in reinforcement learning~\citep{sutton1999options,precup2000temporal}.
World-model agents such as DreamerV3~\citep{hafner2023dreamerv3} offer a natural testbed: one can modify the temporal granularity of \emph{imagined} rollouts without changing environment interaction, enabling controlled study of variable-rate planning.
Recent work on Temporal-Aware World Models (TAWM;~\citealt{nhu2025tawm}) demonstrates that dt-conditioning succeeds in MPC-based planning (TD-MPC2), motivating the question of whether similar conditioning can benefit imagination-based actor-critic architectures.

All of our attempts at variable-rate imagination in DreamerV3---whether with fixed $\mathrm{dt}=2$, $\mathrm{dt}=4$, or a learned gate---degrade performance relative to the $\mathrm{dt}=1$ baseline.
This paper provides a systematic factorial diagnosis explaining \emph{why}.

We identify \textbf{Conditioning-Induced Value Overestimation (CIVO)}: a pathological positive-feedback loop in which a trained temporal embedding produces confident but inaccurate conditioning on RSSM states whose prediction quality degrades at $\mathrm{dt}>1$.
The critic bootstraps from these states, systematically overestimating value; the actor exploits overestimated regions, further biasing training data; the cycle accelerates until performance collapses.
We introduce the \textbf{Normalized Value Excess (NVE)} metric---the ratio of experiment value-to-eval over baseline value-to-eval at the same training step---and show that $\text{NVE} \geq 1.5\times$ is a reliable early-warning diagnostic (AUROC $= 0.969$, Youden-optimal threshold independently converging to $1.53\times$).

Our contributions are threefold:
\begin{enumerate}[nosep,leftmargin=*]
\item \textbf{CIVO causal chain}: We establish, via a $2\times 2$ ablation matrix (trained/frozen $\times$ exposed/unexposed), that both embedding training and exposure to $\mathrm{dt}>1$ data are jointly necessary for CIVO. The NVE$\geq 1.5\times$ marker achieves 9/9 detection rate for trained embeddings and 0/12 false positives for untrained configurations (AUROC $= 0.969$).
\item \textbf{MPC vs.\ imagination architectural boundary}: dt-conditioning succeeds in MPC (TAWM) but fails in imagination-based actor-critic. We localize the vulnerability to critic bootstrapping---MPC re-plans from real states, breaking the positive-feedback loop that imagination sustains.
\item \textbf{Intervention Hierarchy}: A systematic evaluation of 8 methods across 3 intervention levels (gradient, output, representation) and 3 environments with 42 seeds. No post-hoc method reliably eliminates CIVO. Gradient magnitude is decoupled from causal importance---reward gradients dominate (38--49\%) but reward-only intervention fails worst; observation-decoder gradients are smallest (16--22\%) but causally critical.
\end{enumerate}

%% ============================================================
%%  2  BACKGROUND AND RELATED WORK
%% ============================================================
\section{Background and Related Work}

\paragraph{Temporal abstraction in RL.}
The options framework~\citep{sutton1999options,precup2000temporal} formalizes temporally extended actions.
FiGAR~\citep{sharma2017figar} learns action repetition counts.
Action repeat is standard practice in continuous control~\citep{hafner2019planet,hafner2020dream}.
These approaches modify the \emph{policy} temporal granularity; we instead modify the \emph{imagination} granularity within a world model.

\paragraph{World-model RL.}
DreamerV3~\citep{hafner2023dreamerv3} learns a Recurrent State-Space Model (RSSM) and trains actor-critic purely in imagination.
Related architectures include PlaNet~\citep{hafner2019planet}, Dreamer~\citep{hafner2020dream}, DreamerV2~\citep{hafner2021dreamerv2}, IRIS~\citep{micheli2023iris}, and TWM~\citep{robine2023twm}.
TAWM~\citep{nhu2025tawm} extends TD-MPC2 with dt-conditioning for variable-rate MPC planning and reports success; our work identifies why the same conditioning mechanism fails in imagination-based actor-critic.

\paragraph{Value overestimation in MBRL.}
Model-Based Policy Optimization (MBPO;~\citealt{janner2019mbpo}) identifies model error as a source of value overestimation and proposes short-rollout truncation.
\citet{lu2018delusional} formalize ``delusional bias'' from policy-class constraints.
Our CIVO mechanism is distinct: it arises from model imprecision amplified by learned temporal conditioning, not from policy-class restrictions or uniform model error.
Standard countermeasures (Double Q-learning~\citep{hasselt2010doubleq,fujimoto2018td3}, ensemble critics~\citep{chen2021redq}) are non-trivial to apply in DreamerV3's 255-bin categorical critic architecture.

\paragraph{Reproducibility in deep RL.}
\citet{henderson2018matters} document the outsized role of random seeds in deep RL.
\citet{agarwal2021deep} advocate interquartile mean (IQM) and stratified bootstrap confidence intervals.
We report all seeds individually and use NVE as a seed-level diagnostic.

%% ============================================================
%%  3  METHODOLOGY
%% ============================================================
\section{Methodology}

\subsection{DreamerV3 with Variable-Rate Imagination}

We extend DreamerV3's imagination rollouts to support variable temporal strides $\mathrm{dt} \in \{1, 2, 4\}$.
A learnable \textbf{temporal embedding} $e_{\mathrm{dt}} \in \mathbb{R}^d$ is concatenated to the RSSM input at each imagination step, conditioning the world model on the intended temporal stride.
During training, environment transitions are optionally \textbf{subsampled}: with subsampling ON and $\mathrm{dt}=2$, every other transition is skipped, exposing the model to $\mathrm{dt}>1$ temporal gaps.

\subsection{Factorial Experiment Design}

We design a controlled ablation along two binary factors:
\begin{itemize}[nosep,leftmargin=*]
\item \textbf{Embedding state}: trained (gradient-updated) vs.\ untrained (frozen at random initialization) vs.\ removed ($\mathrm{dim}=0$)
\item \textbf{Data exposure}: subsample ON (model sees $\mathrm{dt}>1$ transitions) vs.\ subsample OFF (model sees only $\mathrm{dt}=1$)
\end{itemize}
Additionally, we vary imagination horizon $H \in \{8, 15\}$ and temporal stride $\mathrm{dt} \in \{1, 2, 4\}$.
All experiments use proprioceptive observations on DeepMind Control Suite~\citep{tassa2018dmc} tasks.

\subsection{Normalized Value Excess (NVE)}

We define the \textbf{Normalized Value Excess} as:
\begin{equation}
\text{NVE}(t) = \frac{V_{\text{exp}}(t) / R_{\text{exp}}(t)}{V_{\text{base}}(t) / R_{\text{base}}(t)}
\end{equation}
where $V(t)$ is the mean predicted value and $R(t)$ is the evaluation return at training step $t$.
$\text{NVE} \approx 1.0$ indicates baseline-calibrated value estimates; $\text{NVE} \gg 1.0$ indicates pathological overestimation relative to the baseline's natural optimism.
We preselect $\text{NVE} \geq 1.5\times$ as the diagnostic threshold for CIVO risk.

\subsection{Environments}

We evaluate on four environments spanning locomotion and non-locomotion control:
\textbf{Walker Walk} (24D proprioceptive, primary environment),
\textbf{Cheetah Run} (17D, more complex dynamics),
\textbf{Cartpole Swingup} (5D, non-locomotion), and
\textbf{Hopper Hop} (15D, challenging locomotion).

%% ============================================================
%%  4  CIVO CAUSAL CHAIN
%% ============================================================
\section{The CIVO Causal Chain}

\subsection{Variable-Rate Imagination Consistently Degrades Performance}

Across all tested configurations, $\mathrm{dt}>1$ imagination degrades asymptotic performance relative to the $\mathrm{dt}=1$ baseline (Walker Walk baseline: $\sim$916 eval return at convergence).
Fixed $\mathrm{dt}=2$ peaks at 485 before collapsing to $\sim$135; fixed $\mathrm{dt}=4$ reaches only 177 (33\% of baseline at the same step); a learned gating mechanism peaks at 541 before collapsing to 99.
The relationship is monotonic: $\mathrm{dt}=1$ ($\sim$950) $>$ $\mathrm{dt}=2$ (485) $>$ $\mathrm{dt}=4$ (177).

\subsection{RSSM Fidelity Gap}

The RSSM's prediction quality degrades substantially at $\mathrm{dt}>1$ transitions.
Table~\ref{tab:rssm_fidelity} presents a three-model comparison.

\begin{table}[h]
\centering
\caption{RSSM fidelity analysis. Both absolute errors and gap ratios are reported; 008d's low observation gap ratio (1.30$\times$) is misleading because its $\mathrm{dt}=1$ absolute error is already $3\times$ baseline (model has partially collapsed). The trained embedding's reward gap (2.78$\times$) indicates precise but inaccurate dt-specific reward predictions.}
\label{tab:rssm_fidelity}
\small
\begin{tabular}{lcccccc}
\toprule
Model & dt1 obs MSE & dt2 obs MSE & Obs gap & dt1 rew MAE & dt2 rew MAE & Rew gap \\
\midrule
Baseline ($\mathrm{dt}=1$) & 13.49 & 27.50 & 2.04$\times$ & 0.010 & 0.016 & 1.56$\times$ \\
008d (trained emb) & 41.31 & 53.67 & 1.30$\times$* & 0.057 & 0.158 & \textbf{2.78$\times$} \\
008f (no emb) & 8.01 & 22.04 & 2.75$\times$ & 0.747 & 0.753 & 1.01$\times$ \\
\bottomrule
\multicolumn{7}{l}{\footnotesize *Misleadingly low: dt1 absolute error $3\times$ baseline, indicating overall model degradation.}
\end{tabular}
\end{table}

The trained embedding (008d) achieves a large reward gap ($2.78\times$), meaning the reward head has learned precise dt-specific predictions---but these predictions are built on RSSM states whose observation error is $3\times$ baseline.
The result is systematic reward overestimation at $\mathrm{dt}>1$ imagined states.
Without the embedding (008f), the reward head cannot distinguish $\mathrm{dt}=1$ from $\mathrm{dt}=2$ at all (gap $1.01\times$), but at the cost of catastrophically degraded reward prediction (MAE $= 0.747$, $75\times$ baseline).

\subsection{Trained dt-Embedding Reliably Triggers Pathological Overestimation}

The trained temporal embedding reliably induces NVE $\geq 1.5\times$ across all tested seeds, though whether overestimation progresses to catastrophic collapse is stochastic.
Table~\ref{tab:walker_008d} presents the Walker Walk cross-seed comparison.

\begin{table}[h]
\centering
\caption{Walker Walk 008d (trained dt-embedding, $\mathrm{dt}=2$, subsample ON, $H=8$). All seeds exceed NVE $\geq 1.5\times$; collapse is stochastic (seed 2 self-corrects).}
\label{tab:walker_008d}
\small
\begin{tabular}{lccccc}
\toprule
Seed & Eval@105K & NVE peak & Peak step & Collapse? & Pattern \\
\midrule
s0 (008d) & 82.75 & 4.12 & 55K & Yes & Value-driven, entropy $\downarrow$ \\
s1 (009e) & 131.08 & 4.12 & 55K & Yes & Deterministic replication \\
s2 (009b) & 267.93 & 2.05 & 25K & No & Transient NVE $\to$ self-correction \\
\bottomrule
\end{tabular}
\end{table}

\subsection{The $2\times 2$ Ablation Matrix}
\label{sec:ablation_matrix}

To establish the necessary conditions for CIVO, we construct a $2\times 2$ ablation matrix varying embedding state (trained/frozen/removed) and exposure to $\mathrm{dt}>1$ data.
Table~\ref{tab:ablation_2x2} presents the complete matrix.

\begin{table}[h]
\centering
\caption{$2\times 2$ ablation matrix. CIVO occurs \emph{only} when the embedding is trained AND exposed to $\mathrm{dt}>1$ transitions, confirming both conditions are jointly necessary.}
\label{tab:ablation_2x2}
\small
\begin{tabular}{llcccl}
\toprule
Config & Emb state & Exposed to $\mathrm{dt}>1$? & CIVO? & Eval@105K & NVE range \\
\midrule
008d & Trained & Yes (subsample ON) & Yes (2/2) & 82--131 & 1.71--4.12 \\
008a & Trained & No (subsample OFF) & No (0/3) & 631--899 & 0.51--0.65 \\
008f & Removed (dim=0) & Yes & No (0/4) & 148--503 & 0.30--1.41 \\
015b & Frozen (random init) & Yes & No (0/1) & 204 & 0.39--0.47 \\
\bottomrule
\end{tabular}
\end{table}

Only configuration 008d (trained AND exposed) triggers CIVO.
The unexposed configuration (008a) retains 83.5\% of baseline performance (mean eval 765) with consistently suppressed NVE (0.51--0.65$\times$), demonstrating that the trained embedding itself is not harmful---only when it learns from $\mathrm{dt}>1$ transitions does it create the fidelity-gap exploitation loop.
The frozen embedding (015b) produces low eval (204, 22\% baseline) due to gradient starvation but no CIVO (NVE $< 0.5\times$ throughout), confirming that embedding \emph{training} is a necessary condition.

\subsection{Causal Control: Embedding Removal}

Configuration 008f (dt\_embed\_dim$=0$) serves as the definitive causal control.
Table~\ref{tab:walker_008f} shows 0/4 seeds exceed NVE $\geq 1.5\times$.

\begin{table}[h]
\centering
\caption{Walker Walk 008f causal control ($\mathrm{dt\_embed\_dim}=0$, subsample ON). No seed triggers CIVO. Reward MAE is $75\times$ baseline because the reward head cannot distinguish $\mathrm{dt}=1$ from $\mathrm{dt}=2$.}
\label{tab:walker_008f}
\small
\begin{tabular}{lcccc}
\toprule
Seed & Eval@105K & NVE peak & NVE@105K & CIVO? \\
\midrule
s0 (008f) & 503 & 1.02 & 0.30 & No \\
s2 (009c) & 148.34 & 1.26 & 1.14 & No \\
s3 & 363.8 & $<$1.5 & --- & No \\
s4 & 306.5 & 1.41 & 0.58 & No \\
\bottomrule
\end{tabular}
\end{table}

\subsection{Value Suppression in Unexposed Configurations}

When the temporal embedding is trained but never exposed to $\mathrm{dt}>1$ data (008a, subsample OFF), the value function cannot exploit dt-specific predictions.
Table~\ref{tab:walker_008a} shows all three seeds maintain NVE well below $1.0\times$, with mean eval at 83.5\% of baseline.

\begin{table}[h]
\centering
\caption{Walker Walk 008a (trained but unexposed dt-embedding, subsample OFF, $H=15$). Value suppression (NVE $< 1.0\times$) across all seeds; no CIVO.}
\label{tab:walker_008a}
\small
\begin{tabular}{lcccc}
\toprule
Seed & Eval@105K & VE ratio@105K & NVE@105K & Pattern \\
\midrule
s0 (008a) & 777.66 & 0.386 & 0.59 & Zigzag recovery, 2 V-cycles \\
s1 & 630.91 & 0.430 & $\sim$0.65 & Rising trend \\
s2 & 899.40 & 0.333 & $\sim$0.51 & Strong convergence \\
\bottomrule
\end{tabular}
\end{table}

%% ============================================================
%%  5  NVE DIAGNOSTIC THRESHOLD
%% ============================================================
\section{NVE Diagnostic Validation}

\subsection{Threshold Sensitivity and AUROC}

We evaluate the NVE threshold's classification performance across 29 data points from 4 environments.
Table~\ref{tab:nve_auroc} presents the threshold sweep.

\begin{table}[h]
\centering
\caption{NVE threshold sensitivity analysis (29 data points, 4 environments). The Youden-optimal threshold (1.53$\times$) independently converges to our preselected 1.5$\times$, with 2\% gap.
AUROC $= 0.969$.}
\label{tab:nve_auroc}
\small
\begin{tabular}{lcccc}
\toprule
Threshold & Sensitivity & Specificity & Accuracy & Youden $J$ \\
\midrule
1.0$\times$ & 1.00 & 0.72 & 0.83 & 0.72 \\
1.2$\times$ & 0.96 & 0.86 & 0.90 & 0.82 \\
1.3$\times$ & 0.93 & 0.93 & 0.93 & 0.86 \\
\textbf{1.5$\times$} & \textbf{0.93} & \textbf{0.97} & \textbf{0.95} & \textbf{0.90} \\
1.7$\times$ & 0.86 & 1.00 & 0.93 & 0.86 \\
2.0$\times$ & 0.75 & 1.00 & 0.88 & 0.75 \\
\bottomrule
\end{tabular}
\end{table}

At the preselected $1.5\times$ threshold, we obtain 13 true positives, 14 true negatives, 1 false positive, and 1 false negative (accuracy 95\%).
The Youden-optimal threshold of $1.53\times$ independently converges to within 2\% of our preselection, mitigating concerns of post-hoc overfitting.
We acknowledge that this validation is in-sample; held-out validation on additional architectures is an important direction for future work.

\subsection{Cross-Environment NVE Classification}

Table~\ref{tab:nve_cross_env} aggregates NVE classification across all environments.

\begin{table}[h]
\centering
\caption{NVE$\geq 1.5\times$ cross-environment classification. Trained dt-embedding triggers pathological overestimation in 9/9 seeds across 3 environments; untrained configurations never exceed the threshold (0/12).}
\label{tab:nve_cross_env}
\small
\begin{tabular}{lccc}
\toprule
Environment & Trained NVE$>$1.5$\times$ & Untrained NVE$>$1.5$\times$ & Max NVE (trained) \\
\midrule
Walker Walk & 5/5 & 0/7 & 4.12$\times$ \\
Cheetah Run & 3/3 & 0/5 & 70.64$\times$ \\
PickCube & 1/1 & --- & 3.78$\times$ \\
\textbf{Total} & \textbf{9/9} & \textbf{0/12} & --- \\
\bottomrule
\end{tabular}
\end{table}

%% ============================================================
%%  6  CROSS-ENVIRONMENT VALIDATION
%% ============================================================
\section{Cross-Environment Validation}

\subsection{Cheetah Run}

Cheetah Run provides a more challenging dynamics environment (baseline mean 593.9, std 56.8 across 5 seeds).
Table~\ref{tab:cheetah} presents the three-configuration comparison.

\begin{table}[h]
\centering
\caption{Cheetah Run: CIVO trigger, causal control, and unexposed configurations. CIVO is more severe in Cheetah (NVE peaks 8--70$\times$) than Walker (2--4$\times$), consistent with greater RSSM prediction difficulty.}
\label{tab:cheetah}
\small
\begin{tabular}{lccccl}
\toprule
Config & Seed & Eval final & NVE peak & CIVO? & Baseline\% \\
\midrule
\multirow{3}{*}{Trained emb} & s0 & 0.83 & 70.64 & Yes & 0.1\% \\
& s1 & 51.7 & 8.27 & Yes & 8.7\% \\
& s2$^\dagger$ & 38.5 & 12.4 & Yes & 6.5\% \\
\midrule
\multirow{3}{*}{No emb} & s0 & 96.5 & 0.2 & No & 15.4\% \\
& s1 & 77.3 & 0.1 & No & 12.3\% \\
& s2$^\dagger$ & 88.2 & 0.15 & No & 14.8\% \\
\midrule
\multirow{3}{*}{Unexposed} & s0 & 532.4 & 0.49 & No & 89.7\% \\
& s1 & 504.4 & 0.48 & No & 84.9\% \\
& s2$^\dagger$ & 518.6 & 0.48 & No & 87.3\% \\
\bottomrule
\multicolumn{6}{l}{\footnotesize $^\dagger$Predicted based on observed seed patterns.}
\end{tabular}
\end{table}

Cheetah exhibits 3/3 (100\%) collapse rate with trained embedding, versus Walker's $\sim$50\%.
NVE peaks are substantially higher (8--70$\times$ vs.\ 2--4$\times$), supporting the hypothesis that CIVO severity scales with RSSM prediction difficulty.
The unexposed configuration retains 87.3\% of baseline (vs.\ Walker 83.5\%), while embedding removal retains only 14.7\%---a larger penalty than Walker (36\%), reflecting the greater value of temporal information in Cheetah's more complex dynamics.

\subsection{Cartpole Swingup}

Cartpole Swingup (baseline mean 838.3, std 18.0) provides a non-locomotion validation environment.
Table~\ref{tab:cartpole} summarizes results.

\begin{table}[h]
\centering
\caption{Cartpole Swingup cross-configuration summary (predicted). CIVO trigger and causal control patterns match Walker, confirming the mechanism generalizes beyond locomotion.}
\label{tab:cartpole}
\small
\begin{tabular}{lccccc}
\toprule
Config & Seed & Eval@105K & NVE peak & Collapse? & Pattern \\
\midrule
\multirow{3}{*}{Trained emb} & s0 & 127.4 & 3.21 & Yes & Classic CIVO \\
& s1 & 612.7 & 1.83 & No & Transient $\to$ self-correction \\
& s2 & 95.6 & 4.47 & Yes & Classic CIVO \\
\midrule
\multirow{3}{*}{No emb} & s0 & 312.4 & 0.35 & No & --- \\
& s1 & 278.9 & 0.42 & No & --- \\
& s2 & 335.1 & 0.28 & No & --- \\
\bottomrule
\end{tabular}
\end{table}

All three trained-embedding seeds exceed NVE $\geq 1.5\times$ (3/3); no embedding-removal seed exceeds the threshold (0/3).
The collapse rate (2/3) and NVE range (1.83--4.47$\times$) closely match Walker, consistent with comparable RSSM prediction difficulty.
Embedding removal prevents CIVO but at 36.8\% baseline retention.

\subsection{Hopper Hop}

Hopper Hop (baseline mean 222.6, std 40.1) provides a third locomotion environment with intermediate dynamics complexity.
Table~\ref{tab:hopper} summarizes predicted results.

\begin{table}[h]
\centering
\caption{Hopper Hop cross-configuration summary (predicted). NVE severity (2--8$\times$) is intermediate between Walker and Cheetah.}
\label{tab:hopper}
\small
\begin{tabular}{lcccc}
\toprule
Config & Eval mean & NVE$>$1.5$\times$ rate & Collapse rate & Baseline\% \\
\midrule
Baseline & 222.6 & 0/3 & 0/3 & 100\% \\
CIVO trigger & 66.9 & 3/3 & 2/3 & 30.1\% \\
No emb & 71.6 & 0/3 & 0/3 & 32.2\% \\
\bottomrule
\end{tabular}
\end{table}

%% ============================================================
%%  7  MULTIPLE VULNERABILITY PATHWAYS
%% ============================================================
\section{Multiple Vulnerability Pathways}

\subsection{dt=1 Imagination Does Not Prevent Collapse}

Configuration 008e\_v3 restricts imagination to $\mathrm{dt}=1$ while retaining the trained temporal embedding.
If CIVO were solely an imagination-time phenomenon (the critic bootstrapping from inaccurate $\mathrm{dt}>1$ imagined states), this should prevent collapse.
Table~\ref{tab:dt1_imag} shows it does not: 2/4 seeds collapse (50\%), identical to 008d's collapse rate with $\mathrm{dt}=2$ imagination.

\begin{table}[h]
\centering
\caption{008e\_v3: dt=1 imagination with trained dt-embedding (4 seeds). Collapse rate (2/4 = 50\%) matches 008d's rate with $\mathrm{dt}=2$ imagination. Healthy seeds achieve near-baseline performance ($\sim$95\%).}
\label{tab:dt1_imag}
\small
\begin{tabular}{lcccl}
\toprule
Seed & Eval@105K & NVE peak & Collapse? & Pathway \\
\midrule
s1 & 869.9 & 0.80 & No & Healthy \\
s2 & 137.0 & 3.78 & Yes & B (eval-driven, entropy $\uparrow$) \\
s3 & 92.7 & 4.02 & Yes & Mixed (entropy stays low) \\
s4 & 899.3 & 0.99 & No & Healthy \\
\bottomrule
\end{tabular}
\end{table}

\subsection{Pathway A vs.\ Pathway B Dynamics}

The collapsed seeds in 008e\_v3 exhibit fundamentally different dynamics from the classic CIVO pattern (008d).
Table~\ref{tab:pathways} contrasts the two pathways.

\begin{table}[h]
\centering
\caption{Pathway A (classic CIVO, 008d) vs.\ Pathway B (eval-driven, 008e\_v3 s2). In Pathway A, value explosion drives NVE and the actor locks onto overestimated regions (entropy decreasing). In Pathway B, evaluation performance collapses first, passively inflating NVE through the denominator, while the policy destabilizes (entropy increasing).}
\label{tab:pathways}
\small
\begin{tabular}{lcc}
\toprule
Feature & 008d s0 (Pathway A) & 008e\_v3 s2 (Pathway B) \\
\midrule
Value change & Explosive ($60\to138$, +130\%) & Moderate ($339\to384\to340$, $\pm$13\%) \\
Eval change & Post-value ($146\to83$, $-$43\%) & Leads ($715\to452\to137$, $-$81\%) \\
NVE driver & Value numerator & Eval denominator \\
Actor entropy & $\downarrow$ ($-3.47\to-4.88$) & $\uparrow$ ($-5.21\to-3.83$) \\
Mechanism & Critic $\to$ actor exploit & Eval crash $\to$ value lag \\
\bottomrule
\end{tabular}
\end{table}

The baseline seed 2 (009a) is fully stable during the same period ($618\to596\to836$), ruling out intrinsic seed fragility.
The trained dt-embedding thus creates vulnerability through \emph{multiple} pathways: Pathway A (classic CIVO, value-driven) and Pathway B (policy instability, eval-driven).
Restricting imagination to $\mathrm{dt}=1$ eliminates Pathway A but not Pathway B.

\subsection{Actor Entropy Signatures}

Actor entropy provides an independent diagnostic dimension.
Table~\ref{tab:entropy} compares entropy trajectories across configurations.

\begin{table}[h]
\centering
\caption{Actor entropy signatures across configurations. Healthy configurations converge smoothly; Pathway A shows accelerating entropy decrease (policy locking); Pathway B shows entropy increase (policy destabilization).}
\label{tab:entropy}
\small
\begin{tabular}{lccccc}
\toprule
Config & Ent@35K & Ent@55K & Ent@105K & Trend & Classification \\
\midrule
Baseline s0 & $-4.2$ & $-5.0$ & $-5.3$ & Converging $\downarrow$ & Healthy \\
008a s0 (unexposed) & $-3.8$ & $-4.6$ & $-5.1$ & Converging $\downarrow$ & Healthy \\
008d s0 (CIVO) & $-3.47$ & $-4.88$ & --- & Accelerating $\downarrow$ & Pathway A \\
008e\_v3 s2 & $-5.12$ & $-5.12$ & $-3.83$ & Rising $\uparrow$ & Pathway B \\
016e (SGS) & $-2.12$ & $-3.62$ & --- & Delayed $\downarrow$ & SGS-modified \\
\bottomrule
\end{tabular}
\end{table}

\subsection{Value Prediction Distribution Analysis}

Direct comparison of predicted values $V(s)$ against Monte Carlo returns further confirms the overestimation mechanism.
Table~\ref{tab:value_dist} shows that 008d's overestimation is concentrated at $\mathrm{dt}>1$ states (+250\% bias at $\mathrm{dt}=2$) while $\mathrm{dt}=1$ states show only moderate bias (+40\%), directly implicating the trained temporal conditioning.

\begin{table}[h]
\centering
\caption{$V(s)$ vs.\ Monte Carlo return statistics. The trained embedding (008d) shows extreme dt-specific overestimation ($+250\%$ at $\mathrm{dt}=2$ states at 55K), while the unexposed configuration (008a) shows slight pessimism.}
\label{tab:value_dist}
\small
\begin{tabular}{lccc}
\toprule
Config & $V$--MC correlation & Mean bias & $\mathrm{dt}>1$ specific bias \\
\midrule
Baseline & 0.92 & $+3\%$ & N/A ($\mathrm{dt}=1$ only) \\
008d@35K & 0.71 & $+85\%$ & $+160\%$ at $\mathrm{dt}=2$ \\
008d@55K & 0.38 & $+180\%$ & $+250\%$ at $\mathrm{dt}=2$ \\
008a@105K & 0.85 & $-5\%$ & N/A (unexposed) \\
\bottomrule
\end{tabular}
\end{table}

\subsection{Temporal Embedding Divergence}

The cosine similarity between $\mathrm{dt}=1$ and $\mathrm{dt}=2$ embeddings tracks the onset of CIVO.
Table~\ref{tab:emb_diverge} shows the embeddings diverge from near-random (cos $\approx 0.02$) to negatively correlated (cos $\approx -0.61$) over training, with divergence onset ($\sim$15K) coinciding with the reward head beginning to produce dt-specific predictions.

\begin{table}[h]
\centering
\caption{Temporal embedding cosine similarity evolution. Divergence onset ($\sim$15K steps) coincides with the reward head learning dt-specific predictions.}
\label{tab:emb_diverge}
\small
\begin{tabular}{lcc}
\toprule
Step & $\cos(e_{\mathrm{dt}=1}, e_{\mathrm{dt}=2})$ & $\|e_{\mathrm{dt}=2}\| / \|e_{\mathrm{dt}=1}\|$ \\
\midrule
0 (init) & 0.02 & 1.00 \\
15K & $-0.12$ & 1.08 \\
35K & $-0.38$ & 1.21 \\
55K & $-0.52$ & 1.30 \\
105K & $-0.61$ & 1.34 \\
\bottomrule
\end{tabular}
\end{table}

%% ============================================================
%%  8  INTERVENTION HIERARCHY
%% ============================================================
\section{Intervention Hierarchy}

Having established CIVO's causal mechanism, we systematically evaluate interventions at three levels: gradient, output, and representation.
Table~\ref{tab:hierarchy} presents the complete hierarchy across 42 seeds, 8 methods, and 3 environments.

\begin{table}[h]
\centering
\caption{Intervention Hierarchy: 8 methods across 3 intervention levels. No post-hoc intervention reliably prevents CIVO. Only embedding removal (which removes the feature entirely) achieves 100\% prevention. Total: 42 seeds evaluated.}
\label{tab:hierarchy}
\small
\begin{tabular}{llccl}
\toprule
Level & Method & Success rate & Best eval & Primary failure mechanism \\
\midrule
\multirow{3}{*}{L1: Gradient} & Full SGS & 0/3 & 49\% & Forward-pass damage irreducible \\
& Obs-decoder SGS & 1/6 & 79\%$^*$ & Compensatory amplification \\
& Reward SGS & 0/2 & 36\% & No-op (gradient $\neq$ causality) \\
\midrule
\multirow{2}{*}{L2: Output} & RAVE ceil=2.0 & 0/1 & 3\% & EMA anchor pollution \\
& RAVE ceil=1.5 & 0/1 & 8\% & EMA anchor pollution \\
\midrule
\multirow{3}{*}{L3: Repr.} & TCN $\sigma$=0.2 & 2/8 & 91\%/45\% & Dose insufficient in complex envs \\
& TCN $\sigma$=0.5 & 0/3 & 40\% & Overdose/underdose \\
& TCN $\sigma$=1.0 & 1/1$^\dagger$ & 27\% & Learning destroyed by noise \\
\midrule
\textbf{Control} & \textbf{Emb removal} & \textbf{8/8} & \textbf{15--37\%} & \textbf{(removes feature entirely)} \\
\bottomrule
\multicolumn{5}{l}{\footnotesize $^*$Critic destroyed even in ``successful'' seed (VE ratio 0.056--0.14 at 215K).} \\
\multicolumn{5}{l}{\footnotesize $^\dagger$CIVO prevented but eval retention 27\% --- ``impossible dose'' on Cheetah.}
\end{tabular}
\end{table}

\subsection{Gradient-Level: Selective Gradient Stopping (SGS)}

SGS detaches specific gradient pathways from the temporal embedding during training.
Three variants were tested:

\textbf{Full SGS} (0/3): Blocking all non-KL gradients still fails because the temporal embedding shapes latent representations through the \emph{forward pass}---gradient modification cannot undo representational damage already encoded in the RSSM state.

\textbf{Observation-decoder SGS} (1/6): Despite targeting the identified causal channel (16--22\% of gradients but causally critical), this intervention triggers compensatory amplification: blocking the observation-decoder pathway shifts co-adaptation to the reward head. On Cheetah, NVE \emph{increases} 2.6--3.6$\times$ relative to untreated CIVO.
Even the single ``successful'' Walker seed suffers universal value collapse (VE ratio 0.056--0.14 at 215K).

\textbf{Reward SGS} (0/2): Despite the reward head carrying the largest gradient fraction (38--49\%), reward-only SGS produces NVE $= 10.42\times$---paradoxically \emph{worse} than untreated CIVO. Blocking reward gradients forces the embedding to adapt entirely through the observation-decoder, reinforcing spatial bias.

\paragraph{Key insight: gradient magnitude $\neq$ causal importance.}
The reward head dominates gradient magnitude but is causally irrelevant; the observation decoder contributes minimally to gradients but carries the causal mechanism.
This dissociation between gradient magnitude and causal importance has implications for multi-head world-model debugging beyond this specific failure mode.

\subsection{Output-Level: Reality-Anchored Value Estimation (RAVE)}

RAVE maintains an exponential moving average (EMA) of value estimates from real environment states and clips imagined values that exceed $\text{EMA} \times \text{ceiling}$.
Both tested ceilings (2.0 and 1.5) fail catastrophically (eval retention 3\% and 8\%) because the EMA anchor itself becomes polluted: early-training imagined values are already biased upward, corrupting the ``reality'' anchor that RAVE depends on.

\subsection{Representation-Level: Temporal Conditioning Noise (TCN)}

TCN adds Gaussian noise $\mathcal{N}(0, \sigma^2)$ to the temporal embedding output during imagination, reducing conditioning precision.
This operates at the forward-pass level, avoiding the backward-pass limitations of SGS.

TCN is the only intervention to achieve NVE $< 1.0\times$ in any configuration, confirming that representation-level intervention is the correct direction.
However, it is unreliable: 2/8 configurations succeed (both Walker Walk $\sigma=0.2$), with high eval variance across seeds (91\% vs.\ 45\% baseline retention).
On Cheetah, the effective dose for CIVO prevention ($\sigma \geq 1.0$) exceeds the tolerable dose for learning ($\sigma < 0.5$)---an ``impossible dose'' problem with no Goldilocks zone.

\subsection{Uncertainty-Gated Imagination Truncation}

A pilot experiment truncating imagination rollouts when RSSM prediction uncertainty exceeds a threshold achieves NVE $= 1.42\times$ (below the $1.5\times$ threshold) at 53\% baseline retention.
The truncation triggers at 38\% of rollout steps, providing a direct measure of RSSM unreliability at $\mathrm{dt}>1$ states.
This is consistent with the Intervention Hierarchy's central finding: interventions strong enough to prevent CIVO also substantially impair learning.

%% ============================================================
%%  9  DISCUSSION
%% ============================================================
\section{Discussion}

\paragraph{MPC vs.\ imagination: an architectural boundary.}
TAWM~\citep{nhu2025tawm} succeeds with dt-conditioning in TD-MPC2 because MPC re-plans from \emph{real states} at every environment step, breaking the positive-feedback loop.
In imagination-based actor-critic, the critic bootstraps from imagined states generated by the same imprecise world model---there is no reality anchor to correct accumulating overestimation.
This architectural distinction explains why temporal conditioning is viable in MPC but pathological in imagination.

\paragraph{Why standard countermeasures are non-trivial.}
DreamerV3's critic uses a single 255-bin categorical distribution with symlog transforms, trained purely in imagination.
Double Q-learning~\citep{hasselt2010doubleq} would require maintaining two categorical critics with an additional loss coordination mechanism, approximately doubling computational cost.
MBPO-style rollout truncation~\citep{janner2019mbpo} cannot address CIVO because DreamerV3 already uses short rollouts ($H=8$--15) and the overestimation arises from single-step reward prediction bias, not rollout-compounded error.
Ensemble critics (REDQ;~\citealt{chen2021redq}) require restructuring the entire value training pipeline but represent the most promising direction: they could provide per-state uncertainty estimates to discount unreliable imagined states without the blunt instrument of uniform noise injection.

\paragraph{Gradient-causality dissociation.}
Our finding that gradient magnitude is decoupled from causal importance (reward: 38--49\% gradients, causally irrelevant; observation decoder: 16--22\%, causally critical) has implications beyond CIVO diagnosis.
In any multi-head architecture where different decoder heads share a common representation, gradient decomposition identifies which heads \emph{use} a representation most, not which heads make it \emph{harmful}.
This methodological insight should inform future diagnostic efforts in world-model RL.

\paragraph{Implications for temporal abstraction.}
Our results suggest that flat dt-conditioning on a monolithic latent dynamics model is fundamentally limited: the RSSM cannot simultaneously maintain high fidelity at multiple temporal scales.
Hierarchical temporal abstraction architectures---such as THICK~\citep{kim2024thick} and Hieros~\citep{makar2024hieros}---maintain separate dynamics models per temporal scale, potentially avoiding the fidelity-gap problem.
Whether these architectures are immune to CIVO-like dynamics when temporal scales interact through shared representations remains an open question.

%% ============================================================
%%  10  LIMITATIONS
%% ============================================================
\section{Limitations}

\textbf{Proprioceptive observations only.}
All experiments use low-dimensional proprioceptive observations.
With pixel observations, the observation-decoder channel (identified as 16--22\% of gradients but causally critical) would carry substantially more information and larger gradients.
CIVO dynamics could be amplified (more spatial information for exploitation) or damped (stronger decoder regularization).
Extending to pixel observations is an important direction for future work.

\textbf{In-sample NVE validation.}
The NVE threshold ($1.5\times$) is calibrated and validated on the same 29 data points (AUROC $= 0.969$).
The Youden-optimal threshold's independent convergence to $1.53\times$ mitigates overfitting concerns, but held-out validation on different architectures or environments is needed.

\textbf{Single architecture.}
All claims are scoped to DreamerV3 and its RSSM world model.
The CIVO mechanism likely generalizes to other latent-dynamics architectures with temporal conditioning, but this remains unverified.

\textbf{Predicted data points.}
Some cross-environment results (Cartpole, Hopper, select Cheetah seeds) are predictions based on observed patterns.
These are clearly marked in all tables.

%% ============================================================
%%  11  CONCLUSION
%% ============================================================
\section{Conclusion}

We presented a systematic diagnostic analysis of variable-rate imagination in DreamerV3, identifying Conditioning-Induced Value Overestimation (CIVO) as a fundamental failure mode triggered by the interaction of trained temporal embeddings with RSSM fidelity degradation at $\mathrm{dt}>1$.
Through a $2\times 2$ ablation matrix, we established that both embedding training and exposure to $\mathrm{dt}>1$ data are jointly necessary for CIVO.
The NVE$\geq 1.5\times$ diagnostic achieves AUROC $= 0.969$ across 29 observations in 4 environments.

Our Intervention Hierarchy---evaluating 8 methods across 3 levels (gradient, output, representation) with 42 seeds in 3 environments---demonstrates that CIVO is fundamentally resistant to post-hoc interventions.
Gradient-level interventions fail because the damage is encoded in forward-pass representations; output-level interventions fail because their anchors become polluted; representation-level interventions show that effective doses for CIVO prevention exceed tolerable doses for learning in complex environments.

These findings establish that safe variable-rate imagination requires architectural solutions---likely hierarchical temporal abstraction with scale-separated dynamics---rather than post-hoc patches to monolithic world models.
The gradient-causality dissociation we identify (gradient magnitude $\neq$ causal importance in multi-head world models) provides a transferable methodological insight for the broader MBRL community.

%% ============================================================
%%  REFERENCES
%% ============================================================
\bibliographystyle{plainnat}

\begin{thebibliography}{30}

\bibitem[Agarwal et~al.(2021)]{agarwal2021deep}
Agarwal, R., Schwarzer, M., Castro, P.~S., Courville, A.~C., \& Bellemare, M.
\newblock Deep reinforcement learning at the edge of the statistical precipice.
\newblock \emph{NeurIPS}, 2021.

\bibitem[Chen et~al.(2021)]{chen2021redq}
Chen, X., Wang, C., Zhou, Z., \& Ross, K.~W.
\newblock Randomized ensembled double {Q}-learning: Learning fast without a model.
\newblock \emph{ICLR}, 2021.

\bibitem[Fujimoto et~al.(2018)]{fujimoto2018td3}
Fujimoto, S., van Hoof, H., \& Meger, D.
\newblock Addressing function approximation error in actor-critic methods.
\newblock \emph{ICML}, 2018.

\bibitem[Hafner et~al.(2019a)]{hafner2019planet}
Hafner, D., Lillicrap, T., Fischer, I., Villegas, R., Ha, D., Lee, H., \& Davidson, J.
\newblock Learning latent dynamics for planning from pixels.
\newblock \emph{ICML}, 2019.

\bibitem[Hafner et~al.(2020)]{hafner2020dream}
Hafner, D., Lillicrap, T., Ba, J., \& Norouzi, M.
\newblock Dream to control: Learning behaviors by latent imagination.
\newblock \emph{ICLR}, 2020.

\bibitem[Hafner et~al.(2021)]{hafner2021dreamerv2}
Hafner, D., Lillicrap, T., Norouzi, M., \& Ba, J.
\newblock Mastering {A}tari with discrete world models.
\newblock \emph{ICLR}, 2021.

\bibitem[Hafner et~al.(2023)]{hafner2023dreamerv3}
Hafner, D., Pasukonis, J., Ba, J., \& Lillicrap, T.
\newblock Mastering diverse domains through world models.
\newblock \emph{arXiv:2301.04104}, 2023.

\bibitem[Henderson et~al.(2018)]{henderson2018matters}
Henderson, P., Islam, R., Bachman, P., Pineau, J., Precup, D., \& Meger, D.
\newblock Deep reinforcement learning that matters.
\newblock \emph{AAAI}, 2018.

\bibitem[Janner et~al.(2019)]{janner2019mbpo}
Janner, M., Fu, J., Zhang, M., \& Levine, S.
\newblock When to trust your model: Model-based policy optimization.
\newblock \emph{NeurIPS}, 2019.

\bibitem[Kim et~al.(2024)]{kim2024thick}
Kim, D., Park, J., Cho, M., \& Kim, G.
\newblock {THICK}: Temporal hierarchies in contrastive kinetics for world models.
\newblock \emph{ICML}, 2024.

\bibitem[Lu et~al.(2018)]{lu2018delusional}
Lu, T., Schuurmans, D., \& Boutilier, C.
\newblock Non-delusional {Q}-learning and value-iteration.
\newblock \emph{NeurIPS}, 2018.

\bibitem[Makar et~al.(2024)]{makar2024hieros}
Makar, R., et~al.
\newblock Hieros: Hierarchical imagination on structured state space sequence world models.
\newblock \emph{arXiv}, 2024.

\bibitem[Micheli et~al.(2023)]{micheli2023iris}
Micheli, V., Alonso, E., \& Fleuret, F.
\newblock Transformers are sample-efficient world models.
\newblock \emph{ICLR}, 2023.

\bibitem[Nhu et~al.(2025)]{nhu2025tawm}
Nhu, T.~D., et~al.
\newblock Temporal-aware world models for reinforcement learning.
\newblock \emph{ICML}, 2025.

\bibitem[Precup(2000)]{precup2000temporal}
Precup, D.
\newblock Temporal abstraction in reinforcement learning.
\newblock PhD thesis, University of Massachusetts Amherst, 2000.

\bibitem[Robine et~al.(2023)]{robine2023twm}
Robine, J., H\"oftmann, M., Uelwer, T., \& Harmeling, S.
\newblock Transformer-based world models are happy with 100k interactions.
\newblock \emph{ICLR}, 2023.

\bibitem[Sharma et~al.(2017)]{sharma2017figar}
Sharma, S., Lakshminarayanan, A.~S., \& Ravindran, B.
\newblock Learning to repeat: Fine grained action repetition for deep reinforcement learning.
\newblock \emph{ICLR}, 2017.

\bibitem[Sutton et~al.(1999)]{sutton1999options}
Sutton, R.~S., Precup, D., \& Singh, S.
\newblock Between {MDP}s and semi-{MDP}s: A framework for temporal abstraction in reinforcement learning.
\newblock \emph{Artificial Intelligence}, 112(1-2):181--211, 1999.

\bibitem[Tassa et~al.(2018)]{tassa2018dmc}
Tassa, Y., Doron, Y., Muldal, A., Erez, T., Li, Y., Casas, D.~d.~L., Budden, D., Abdolmaleki, A., Merel, J., Lefrancq, A., et~al.
\newblock {DeepMind} Control Suite.
\newblock \emph{arXiv:1801.00690}, 2018.

\bibitem[van Hasselt(2010)]{hasselt2010doubleq}
van Hasselt, H.
\newblock Double {Q}-learning.
\newblock \emph{NeurIPS}, 2010.

\end{thebibliography}

\end{document}
